\chapter{Discussion and Critical Analysis}
\label{chap:discussion}

\section{The Main Lesson: No Backend Wins Everywhere}

The final result of this project is best understood as an adaptive-heuristic result. Backend performance depends on workload shape, checkpoint layout, and integration layer, which means checkpoint loading should be driven by workload-aware selection rather than by one fixed backend.

This is the most important correction to the early narrative. The project began with the hope that a single backend --- ideally \texttt{io\_uring} on Linux --- might dominate across the board. The final evidence does not support that conclusion. At the Rust layer, \texttt{SafeTensors} and \texttt{ServerlessLLM} prefer different backends in different size regimes. At the Python layer, binding overhead changes the practical ordering. At the vLLM layer, model initialisation dominates while steady-state decode is largely insensitive to the loading backend.

That is why the final system is organised around an adaptive \texttt{default} policy. In the clearest high-level form, the heuristic is sync for smaller eager loads and io\_uring for larger ones, with format-specific refinements where the access pattern changes.

\section{Why the SafeTensors and ServerlessLLM Stories Diverge}

The results diverge because the two formats induce different I/O structures.

\textbf{SafeTensors} is fundamentally a whole-file or multi-shard eager-loading problem. The dominant costs are reading large contiguous files and coordinating shard-level parallelism. In this regime, \texttt{sync} remains highly competitive for smaller models because its thread-parallel shard loading is simple and effective. As shard count and checkpoint size increase, the multi-worker \texttt{io\_uring} backend becomes increasingly attractive because it can sustain more in-flight work without the same blocking behaviour.

\textbf{ServerlessLLM} is fundamentally a range-heavy partitioned-loading problem. The dominant costs are request grouping, range batching, and fragmented access across partitions. In this regime, the synchronous backend is rarely ideal. The competition is instead between a well-grouped async path and the multi-worker \texttt{io\_uring} path. This is why the explicit \texttt{async} backend is strong for some mid-sized \texttt{ServerlessLLM} cases even when it is weak for \texttt{SafeTensors}.

The practical implication is clear: backend choice cannot be separated from layout choice. A loader that looks strong on \texttt{SafeTensors} does not automatically remain strong on a partitioned format, so any good heuristic has to look at workload structure rather than model size alone.

\section{The Architectural Turnaround Matters More Than the Backend Label}

One of the most useful lessons from the project is that backend family names can be misleading if the architecture behind them is poor.

The early single-ring \texttt{io\_uring} backend performed badly enough that it would have been easy to conclude that \texttt{io\_uring} was simply the wrong fit for checkpoint loading. That conclusion would have been wrong. Profiling showed that the prototype underutilised the machine and spent too much time in the single submission/completion path. The problem was not ``Linux async I/O is bad for this workload''; the problem was ``this specific \texttt{io\_uring} architecture is bad for this workload''.

After the redesign to multiple worker-owned rings, the \texttt{io\_uring} backend changed from a negative result into a strong explicit contender. That shift is central to the interpretation of the work. It shows that backend evaluation for checkpoint loading must account not only for API family but also for execution topology, queue ownership, and batching structure.

This also explains why the report should preserve the turnaround story rather than pretending the earlier negative result never happened. The failed prototype was informative because it revealed which aspects of the architecture actually mattered.

\section{Why Python Needed Its Own Optimisation Work}

A naïve reading of the Rust results might suggest that the Python layer should simply inherit the same backend ordering. The experiments show that this is not true. Python introduces substantial extra work: runtime orchestration, tensor conversion, Torch dtype resolution, and dictionary materialisation. Those costs can swamp part of the raw I/O gain.

This is why the binding-specific optimisation pass was necessary. Reusing a shared Tokio runtime and reducing repeated Torch conversion overhead materially improved the public \texttt{default} path on Qwen3-8B. The important engineering lesson is that a fast core does not automatically imply a fast binding layer. If the goal is practical usability, the binding path must be profiled and optimised in its own right.

This does not weaken the backend conclusions. It refines them. The Rust layer answers which backend is best for raw checkpoint loading. The Python layer answers how much of that gain survives crossing the language boundary. The final system needs both answers.

\section{What the vLLM Results Add}

The vLLM results are valuable because they show how backend effects survive contact with a real serving stack.

The first lesson is that backend gains matter most for model initialisation and TTFT. That is where weight loading is still part of the critical path. The second lesson is that decode throughput is much less sensitive to backend choice, because once the model is resident, the GPU compute path dominates.

The third lesson is methodological. The initial failures of the 14B and 32B vLLM runs looked, at first glance, like large-model infeasibility. In fact, they were caused by conservative vLLM harness settings for memory budgeting and maximum model length. Once those settings were corrected, the full matrix completed successfully. This is a useful reminder that end-to-end integration experiments can fail for harness reasons that are orthogonal to the backend under study.

In practical terms, the vLLM matrix strengthens the project's central argument. A useful loading system is not judged only by backend microbenchmarks; it is judged by whether the selected backend helps the actual serving stack where it matters, especially during cold start.

\section{What Matters for Practitioners}

Three practitioner-oriented conclusions follow directly from the completed experiments.

\textbf{First, prefer policy over dogma.} Engineers should resist the temptation to standardise on one backend globally. The results do not justify such a rule. The right practical design is an adaptive selector with explicit override paths.

\textbf{Second, optimise at the right layer.} If the problem is Rust cold-load latency, backend architecture matters most. If the problem is Python user-facing load time, binding overhead can be equally important. If the problem is serving cold-start latency, vLLM setup and KV-cache configuration must also be considered.

\textbf{Third, evaluate formats and loaders together.} The best eager-loading backend is not a property of the backend alone. It emerges from the interaction between checkpoint layout and loader implementation.

\section{Limitations}

The report's conclusions should be interpreted with several limitations in mind.

\textbf{Hardware specificity.} All final measurements come from one H100-based Linux environment. The ordering of results is likely to generalise more robustly than the absolute timings, but further cross-hardware validation would strengthen that claim.

\textbf{Heuristic imperfection.} The final \texttt{default} selector is much better than the earlier cost-model approach, but it is still not perfect on every Rust matrix row. The medium-size crossover region remains the most difficult area to classify cleanly.

\textbf{Serving-stack dependence.} vLLM results are sensitive to engine settings such as memory budget and maximum model length. This is an unavoidable property of evaluating inside a real serving stack, but it means that the exact numbers are partly harness-dependent.

\textbf{Scope.} The report covers two checkpoint layouts and a six-model ladder. That is much stronger than the earlier prototype setup, but it still does not exhaust the design space of checkpoint formats or storage hardware.

\section{Final Assessment}

The project succeeds if judged against its final, more realistic goal. It does not identify one universally superior checkpoint-loading backend. Instead, it demonstrates that:

\begin{enumerate}
\item backend ranking depends on workload structure,
\item backend architecture matters as much as backend family,
\item Python and serving layers can materially change how backend gains appear in practice, and
\item an adaptive default policy is the most defensible design target.
\end{enumerate}

That final conclusion is stronger than the original, simpler hypothesis would have been. It is also more useful. A practical checkpoint-loading system should not ask users to know which backend is best in every situation. It should use an adaptive heuristic that defaults to the simpler synchronous path when the checkpoint is small enough, moves to io\_uring when larger checkpoints justify the extra machinery, and still leaves room for explicit overrides where experimentation or deployment constraints require them.
